% Generated draft: requires mathematical and source review.
\documentclass[11pt]{article}
\usepackage[margin=25mm]{geometry}
\usepackage{fontspec}
\setmainfont{Latin Modern Roman}
\usepackage{amsmath,amssymb,mathtools,unicode-math}
\setmathfont{Latin Modern Math}
\usepackage{enumitem,needspace,xurl,hyperref}
\setlength{\parindent}{0pt}
\setlength{\parskip}{5pt}
\setlength{\emergencystretch}{3em}
\newcommand{\N}{\mathbb{N}}
\newcommand{\Z}{\mathbb{Z}}
\newcommand{\Q}{\mathbb{Q}}
\newcommand{\R}{\mathbb{R}}
\newcommand{\C}{\mathbb{C}}
\newcommand{\F}{\mathbb{F}}
\newcommand{\A}{\mathbb{A}}
\newcommand{\PP}{\mathbb{P}}
\newcommand{\E}{\mathbb{E}}
\newcommand{\eps}{\varepsilon}
\newcommand{\dd}{\,\mathrm{d}}
\begin{document}
\section*{UnsolvedMath ID: 700692}
\textbf{Source problem code:} catalogue entry
\subsection{Definitions and mathematical statement}
A group \(G\) is periodic if every element of \(G\) has finite order. It is locally finite if every subgroup generated by finitely many elements is finite. For a subset \(S\subseteq G\), define
\[
C_G(S)=\{g\in G:gs=sg\text{ for every }s\in S\},
\qquad C_G(r)=C_G(\{r\}).
\]
An involution is an element of order \(2\). An elementary abelian subgroup of order \(4\) is a subgroup \(R\cong C_2\times C_2\), where \(C_2\) is the cyclic group of order \(2\); its three nonidentity elements are involutions.

For a prime \(p\), the quasicyclic group is
\[
C_{p^\infty}=\{z\in\C^\times:z^{p^n}=1\text{ for some integer }n\geq 0\}.
\]
A group \(H\) is a Chernikov group if it has a normal subgroup of finite index isomorphic to
\[
\prod_{i=1}^{k}C_{p_i^\infty},
\]
where \(k\geq 0\) is finite and each \(p_i\) is prime, with repetitions permitted. The case \(k=0\) includes finite groups.

Let \(G\) be any periodic group containing an elementary abelian subgroup \(R\) of order \(4\). The following are two separate questions.

\textbf{(a)} Does
\[
|C_G(R)|<\infty
\quad\Longrightarrow\quad
G\text{ is locally finite}
\]
hold for every such pair \((G,R)\)?

\textbf{(b)} Does
\[
\bigl(\forall r\in R\setminus\{1\},\ C_G(r)\text{ is a Chernikov group}\bigr)
\quad\Longrightarrow\quad
G\text{ is locally finite}
\]
hold for every such pair \((G,R)\)? No residual finiteness assumption is imposed in either general question.

\textit{Supplied historical qualification (remark of 1999).} The source reports that P. V. Shumyatsky gave a positive answer to part (a) when \(G\) is residually finite, citing the 1998 publication listed in Sources. Here residually finite means that for every \(g\in G\setminus\{1\}\), there exist a finite group \(F\) and a homomorphism \(\varphi:G\to F\) with \(\varphi(g)\neq 1\). This special-case result is recorded as supplied source information and has not been independently verified.
\subsection{Short English statement}
V. P. Shunkov asks whether a periodic group containing a subgroup \(R\cong C_2\times C_2\) must be locally finite under either of two separate hypotheses: the common centralizer of \(R\) is finite, or the centralizer of each of its three involutions is Chernikov. A supplied remark of 1999 reports a positive answer to part (a) for residually finite groups. The registry labels the problem open; that label is supplied curation, not an independently verified assessment.
\subsection{Sources}
\begin{itemize}
\item Supplied local UnsolvedMath registry record: permanent ID \(700692\), exact source problem code \texttt{catalogue entry}, title \emph{Kourovka Notebook Problem 10.77}, upstream revision \texttt{d371ee917ac9647fbb06cd6fc529deec6f3dda82}. The record supplies the status \texttt{open}.
\item \emph{The Kourovka Notebook (Unsolved Problems in Group Theory)}, problem 10.77, proposed by V. P. Shunkov in the 10th issue (1986). The supplied background identifies the 21st-issue edition (2026), arXiv:1401.0300v48. The collection description names E. I. Khukhro and V. D. Mazurov as editors.
\item P. V. Shumyatsky, \emph{Quart. J. Math. Oxford (2)} \textbf{49}, no. 196 (1998), 491--499. The supplied remark of 1999 attributes to this publication a positive answer to part (a) under residual finiteness. No article title, DOI, or URL is supplied.
\end{itemize}
\end{document}
